% --- Section 2: Phase 1 Specifics ---
\section{Classification}
\begin{frame}{Classification}
    The classification task focus on identifying the elemental composition by analyzing spectral signatures.

    \vspace{0.4cm}
    \begin{block}{Analytical stages}
        \begin{itemize}
            \item \textbf{Multi-label classification} across a pool of 41 potential elements.
            \item \textbf{Multi-class classification} across 5 classes defined by the number of elements present in the sample.
        \end{itemize}
    \end{block}

    \vspace{0.3cm}
\end{frame}



\section{Data Generation}
\begin{frame}{Synthetic Data Generation}
    To train models, we developed a Generator that simulates physical variations in XRF spectra.
    \vspace{0.3cm}

    \begin{itemize}
        \item \textbf{Composition:} Samples contain \textbf{0 to 5} elements randomly selected from \textbf{41} possible candidates.
        \item \textbf{Stochasticity:} To ensure model robustness, each spectrum is generated by randomly varying operating conditions and noise levels simulating real-world measurement variability.
    \end{itemize}
\end{frame}



\begin{frame}{Preprocessing Strategies and Architectures}
    \begin{block}{Preprocessing Scaling}
        Tested to handle high dynamic range in spectral counts:
        \begin{itemize}
            \item \textbf{Standard:} 0 mean and 1 variance.
            \item \textbf{MinMax:} Rescaling to $[0, 1]$ range.
            \item \textbf{Log-MinMax:} Logarithmic transformation to compress peaks.
        \end{itemize}
    \end{block}

    \begin{block}{Model Architectures}
        \textbf{Machine Learning:}
        \begin{itemize}
            \item Random Forest, AdaBoost, XGBoost 
        \end{itemize}
        \vspace{0.2cm}
        \textbf{Deep Learning:}
        \begin{itemize}
            \item MLP, CNN, Transformer
        \end{itemize}
    \end{block}
\end{frame}

